\documentclass[10pt,aspectratio=169]{beamer}
%\documentclass[12pt,aspectratio=169]{beamer}

%\mode<presentation>

\usepackage{natbib}
\usepackage{bibentry}
\bibliographystyle{plain}

\usepackage{graphicx}
\usepackage{breqn}
\usepackage{bm}
\usepackage{color}
\usepackage{tabularx}

\usepackage{listings}
\lstset{language=C++}

\usepackage{mathtools}
\DeclarePairedDelimiter\ceil{\lceil}{\rceil}
\DeclarePairedDelimiter\floor{\lfloor}{\rfloor}

\usepackage{animate}

\definecolor{radeonruby}{rgb}{0.87, 0, 0.19}


\usepackage{listings}



%\usepackage[utf8]{inputenc}
\definecolor{bluekeywords}{rgb}{0.63, 0.79, 0.95}
\definecolor{greencomments}{rgb}{0,0.5,0}
\definecolor{redstrings}{rgb}{0.64,0.08,0.08}
\definecolor{xmlcomments}{rgb}{0.5,0.5,0.5}
\definecolor{types}{rgb}{0.17,0.57,0.68}
\lstset{language=C++,
captionpos=b,
%numbers=left,
%numberstyle=\tiny,
frame=single,
showspaces=false,
showtabs=false,
breaklines=true,
showstringspaces=false,
breakatwhitespace=true,
escapeinside={(*@}{@*)},
commentstyle=\color{greencomments},
morekeywords={partial, var, value, get, set},
keywordstyle=\color{bluekeywords},
stringstyle=\color{redstrings},
basicstyle=\ttfamily\small,
}

\usepackage[yyyymmdd]{datetime}
\renewcommand{\dateseparator}{--}

%\usepackage[inline]{asymptote}
\usepackage{asymptote}
\def\asylatexdir{}
\def\asydir{}

\renewcommand{\thefootnote}{\fnsymbol{footnote}}

\def\clap#1{\hbox to 0pt{\hss#1\hss}}
\def\mathllap{\mathpalette\mathllapinternal}
\def\mathrlap{\mathpalette\mathrlapinternal}
\def\mathclap{\mathpalette\mathclapinternal}
\def\mathllapinternal#1#2{\llap{$\mathsurround=0pt#1{#2}$}}
\def\mathrlapinternal#1#2{\rlap{$\mathsurround=0pt#1{#2}$}}
\def\mathclapinternal#1#2{\clap{$\mathsurround=0pt#1{#2}$}}

\def\p{\pause}

\def\({\left(}
\def\){\right)}
\def\[{\left[}
\def\]{\right]}

\def\no#1{\nobreak{#1}} % stop equation breaking in dmath

\begin{asydef}
  // Global Asymptote definitions can be put here.
  import three;
  usepackage("bm");
  texpreamble("\def\V#1{\bm{#1}}");
  // One can globally override the default toolbar settings here:
  // settings.toolbar=true;
\end{asydef}

%\useinnertheme{rectangles} %specifies itemize points to be rectangles
\useoutertheme{default}

\definecolor{darkgreen}{RGB}{0, 100, 0}
\definecolor{brightpurple}{RGB}{160,32,240}

% fg specifies foreground color
% bg specifies background color
% the notation color!x!white mixes the color with 1-x amount of white
%\setbeamercolor{frametitle}{fg=darkgreen!90!black, bg=darkgreen!10!white}
\setbeamercolor{normal text}{fg=white,bg=black}
% defines color of structure elements, e.g., color of text in the
% outline, block titles many more definitions possible. Each element
% predefined by the theme can be overwritten individually

\setbeamerfont{author}{size=\small,series=\bfseries,parent=structure}

\setbeamercolor{background canvas}{bg=white}
%% \usebackgroundtemplate
%% {
%%   \includegraphics[width=\paperwidth,height=\paperheight]{gpu_background_black.png}
%% }
\pgfdeclareimage[width=\paperwidth]{mybackground}{../../../templates/texslides/title.png}
\setbeamertemplate{title page}{
  \begin{picture}(0,0)
    \put(-30,-163){%
      \pgfuseimage{mybackground}
    }
    \put(250,-70){%
      \begin{minipage}[b][45mm][t]{226mm}
        \usebeamerfont{title}{\inserttitle\par}
        \ifx\insertsubtitle\@empty%
        \else%
        {\vskip0.25em%
        \usebeamerfont{subtitle}\usebeamercolor[fg]{subtitle}\insertsubtitle\par}%
        \fi%
        \vskip1cm%              
        \usebeamerfont{author}{\insertauthor\par}
        \usebeamerfont{date}\url{malcolm.roberts@amd.com}\par
        \usebeamerfont{date}\insertdate
      \end{minipage}
    }
  \end{picture}
}

%------------------------------------------------------------------------
%------------------------------------------------------------------------



\def\thetitle{rocFFT MPI}
\def\thesubtitle{HPC Apps Knowledge Sync}
\def\theauthor{rocFFT group}
\def\theauthorinstitute{AMD}

\def\presdate{2025-05-02}

\hypersetup{
  pdftitle=\thetitle{} \thesubtitle{},
  pdfauthor=\theauthor{}
}

\title{\thetitle{}}
\subtitle{\thesubtitle{}}
\author{\theauthor{}}
\institute{\theauthorinstitute}
\date{\presdate}

\setbeamertemplate{frametitle}[default][center]
%\setbeamertemplate{footline}[frame number]
\setbeamertemplate{navigation symbols}{} 

\begin{document}

\nobibliography{../../../refs.bib}

\setbeamertemplate{headline}{}
%\setbeamertemplate{footline}{}

{
  \setbeamercolor{background canvas}{bg=black}
  \setbeamercolor{normal text}{fg=white,bg=black}
  \setbeamercolor{frametitle}{fg=white}
  \setbeamercolor{title}{fg=white}
  \setbeamercolor{alerted text}{fg=white}
  \setbeamercolor{structure}{fg=white}
  \usebeamercolor[fg]{normal text}
  \begin{frame}
    \titlepage
  \end{frame}
}


\setbeamercolor{normal text}{fg=black,bg=white}
\usebeamercolor[fg]{normal text}

\newcolumntype{R}{>{\raggedleft\arraybackslash}X}
\setbeamertemplate{footline}{
   \leavevmode%
   \hbox{%
     \begin{beamercolorbox}[wd=\paperwidth,ht=1ex,dp=1.125ex]{author in head/foot}%
  \begin{tabularx}{\paperwidth}{lR}
    \hfill\insertframenumber/\inserttotalframenumber\hfill\phantom{.}%
    & %
    
\includegraphics[width=0.05\textwidth]{../../../templates/texslides/footerlogo} \textcolor{black}{Together We Advance\begin{animateinline}[autoplay,loop]{1}{\_}\newframe\end{animateinline}}%
    \\%
  \end{tabularx}%
  \end{beamercolorbox}%
  }%
}




\begin{frame}
  \frametitle{\thetitle{}}

  \begin{center}
    \thesubtitle{}
  \end{center}

  \bigskip

  \begin{itemize}
  \item rocFFT multi-GPU/multi-proc API
  \item Communication Library Support
  \item Current Status
  \item Preliminary Scaling Results
  \item Current and Future Work
  \end{itemize}
  
\end{frame}

\begin{frame}
  \frametitle{rocFFT multi-GPU/multi-node API}

  \begin{center}
    Same API for single-proc/multi-proc
  \end{center}

  \bigskip
  
  \begin{itemize}
  \item Divide the input space into \lstinline{rocfft_brick}s using
    \lstinline{rocfft_brick_create}
    \begin{itemize}
    \item Bricks are explicitly associated to devices
    \item Bricks are implicitly associated to an MPI rank
    \end{itemize}
  \item Add the bricks to a \lstinline{rocfft_field} using
    \lstinline{rocfft_field_add_brick}
  \item Add the field to a \lstinline{rocfft_plan_description} using
    \lstinline{rocfft_plan_description_add_infield} or
    \lstinline{rocfft_plan_description_add_outfield}
  \item Add the description to the plan as usual and execute the plan.
  \end{itemize}
  
  \bigskip

  \begin{center}
    Single-GPU/multi-gpu/multi-process transforms use the same library.
  \end{center}
  
\end{frame}


\begin{frame}
  \frametitle{cuFFT API Comparison}


  \begin{itemize}
  \item Single-proc multi-gpu
    \begin{itemize}
    \item cuFFTXT interface.
    \item Only supports \lstinline{CUFFT_XT_FORMAT_INPLACE}
    \end{itemize}
  \item Multi-prod multi-gpu
    \begin{itemize}
    \item cuFFTMP inteface.
    \item Built-in slabs: \lstinline{CUFFT_XT_FORMAT_INPLACE}
    \item Custom brick decomposition: \lstinline{CUFFT_XT_FORMAT_DISTRIBUTED_INPUT}
    \item real/complex transforms only support
      \lstinline{CUFFT_XT_FORMAT_DISTRIBUTED_INPUT}, which means that
      you get to set blocks via
      \lstinline{cufftMpMakePlanDecomposition} allows for a brick
      decomposition.
    \item It seems that only one buffer per rank is allowed.
    \item Makes use of nvshmem
    \end{itemize}
    \item NB: New NVIDIA ARM-based NVPL includes FFT in early access release.
  \end{itemize}

  FIXME: clean this up.
  
\end{frame}

\begin{frame}
  \frametitle{Communication Library Support}

  Compiling rocFFT should work with any GPU-aware MPI library.

  \medskip
  
  Tested with OpenMPI and Cray-MPI

  \bigskip
  
  Release plan:
  \begin{itemize}
  \item Source is always available:
    \begin{itemize}
    \item cmake option \lstinline{ROCFFT_MPI_ENABLE}
      \begin{itemize}
        \item \lstinline{ROCFFT_CRAY_MPI_ENABLE} for cray-mpi to link
          \lstinline{rtl}
      \end{itemize}
    \end{itemize}
  \item Planned OS packages: rocfft, rocfft-openmpi, rocfft-mpich
    \begin{itemize}
    \item rocFFT packages can only support MPI libraries libraries
      which also have OS packages
    \end{itemize}
  \end{itemize}
  \medskip

  
\end{frame}

\begin{frame}
  \frametitle{Current Status}

  \begin{center}
    Source-only early release
  \end{center}
  \begin{itemize}
  \item Initial implementation in 6.3
  \item Performance improvements in 6.4:
    \begin{itemize}
    \item Pencil/slab decomposition using \lstinline{MPI_Ialltoallv}
    \end{itemize}
  \item Upcoming changes in 6.5 or 7.0:
    \begin{itemize}
    \item Use \lstinline{MPI_Ialltoall} instead of
      \lstinline{MPI_Ialltoallv} when possible
    \end{itemize}
  \end{itemize}
  
  \bigskip
  
  Algorithms:
  \begin{itemize}
  \item Gather/Scatter
  \item 1D regular decomposition (slab)
  \item 2D regular decomposition (pencil)
    \begin{itemize}
    \item Performance verification in progress
    \end{itemize}
  \end{itemize}
  
\end{frame}
  
\begin{frame}
  \frametitle{Preliminary Scaling Results}
  \begin{center}
    \includegraphics{m1024.pdf}
    % asy -f pdf datagraphs.asy -u'legendlist="single,double";filenames="1024_3_sp.dat,1024_3_dp.dat"' -o m1024.pdf
    Execution time for a $1024^3$ complex FFT on Frontier with slab
    decomposition
  \end{center}
\end{frame}

\begin{frame}
  \frametitle{Preliminary Scaling Results}
  \begin{center}
    \includegraphics{m2048.pdf}
    % asy -f pdf datagraphs.asy -u'legendlist="single,double";filenames="2048_3_sp.dat,2048_3_dp.dat"' -o m2048.pdf
    Execution time for a $2048^3$ complex FFT on Frontier with slab
    decomposition
  \end{center}
\end{frame}

\begin{frame}
  \frametitle{Preliminary Scaling Results}
  \begin{center}
    \includegraphics{m4096.pdf}
    % asy -f pdf datagraphs.asy -u'legendlist="single,double";filenames="4096_3_sp.dat,4096_3_dp.dat"' -o m4096.pdf
    Execution time for a $4096^3$ complex FFT on Frontier with slab
    decomposition
  \end{center}
\end{frame}

\begin{frame}
  \frametitle{Current and Future Work}

  \begin{itemize}
  \item Ramping-up CI
  \item Scaling up accuracy and performance tests
  \item Improved algorithmic performance
    \begin{itemize}
    \item \lstinline{alltoall} instead of \lstinline{alltoallv}
    \item Optimizing multi-proc with multiple GPUs per rank
    \item Stream- and Kernel-triggered communication
    \item Overcoming latency
    \end{itemize}
  \item First binary release after ROCm-enable OpenMPI is release
  \end{itemize}
  
\end{frame}

\end{document}


% LocalWords:  bg fg frametitle footline lR rocFFT proc rocfft MPI
% LocalWords:  OpenMPI cmake openmpi mpich GROMACS mpi backend FFT
% LocalWords:  alltoall alltoallv ROCm rtl
